% BANKS-SOURCE: pdf=28 print=18 kind=solution-section id=chapter03-implicit-solutions
\begin{bankssolutions}{Solutions to Implicit Exercises for Chapter 3}

% BANKS-SOURCE: pdf=28 print=18 kind=solution id=I-CH03-001
\BanksImplicitSolution{I-CH03-001}
Write $\phi=\phi^{(+)}+\phi^{(-)}$, where $\phi^{(+)}$ contains $a$ and
$\phi^{(-)}$ contains $a^\dagger$.  The vacuum relations
\(\phi^{(+)}|0\rangle=0\) and
\(\langle0|\phi^{(-)}=0\) show that the only two-field vacuum contraction is
\begin{align*}
 C(x,y)&:=\langle0|T\phi(x)\phi(y)|0\rangle\\
 &=\theta(x^0-y^0)
   [\phi^{(+)}(x),\phi^{(-)}(y)]\\
 &\quad+\theta(y^0-x^0)
   [\phi^{(+)}(y),\phi^{(-)}(x)]\\
 &=D_F(x-y).
\end{align*}
Explicitly commuting each annihilation part to the right gives the second
line.  Repeating this commutation proves Wick's operator identity.  Its vacuum expectation vanishes
for an odd number of fields.  For $2r$ fields it is the sum over all complete
pairings,
\[
 \langle0|T\phi(x_1)\cdots\phi(x_{2r})|0\rangle
 =\sum_{P}\prod_{(a,b)\in P}D_F(x_a-x_b).
\]
There are
\[
 N_{2r}=(2r-1)(2r-3)\cdots1
       =\frac{(2r)!}{2^r r!}
\]
such pairings.

Expanding the source exponential now gives
\begin{align*}
 Z_0[J]
 &=\sum_{r=0}^{\infty}\frac{\ii^{2r}}{(2r)!}
   \int\prod_{a=1}^{2r}\dd^D x_a\,J(x_a)
   \sum_P\prod_{(a,b)\in P}D_F(x_a-x_b)\\
 &=\sum_{r=0}^{\infty}\frac{1}{r!}
   \left[-\frac12\int\dd^D x\,\dd^D y\,
    J(x)D_F(x-y)J(y)\right]^r\\
 &=\exp\!\left[-\frac12\int\dd^D x\,\dd^D y\,
    J(x)D_F(x-y)J(y)\right].
\end{align*}
The second equality uses symmetry of the source integrations, which makes all
pairings equal after relabeling.  The zeroth term gives $Z_0[0]=1$.  Finally,
\[
 \left.\frac{\delta^2Z_0}
 {\ii\delta J(x)\,\ii\delta J(y)}\right|_{J=0}
 =D_F(x-y),
\]
so the sign and the factor $1/2$ pass the defining two-point check.

% BANKS-SOURCE: pdf=29 print=19 kind=solution id=I-CH03-002
\BanksImplicitSolution{I-CH03-002}
The time-ordered exponential means the symmetric Dyson series
\[
 Z[J]=\sum_{N=0}^{\infty}\frac{\ii^N}{N!}
 \int\prod_{a=1}^N\dd^D y_a\,J(y_a)
 \langle0|T\phi(y_1)\cdots\phi(y_N)|0\rangle.
\]
A derivative with respect to $J(x)$ can hit any of the $N$ source factors.
The integrand is symmetric under simultaneous exchange of $y_a$ and its
field, so the $N$ terms are equal.  Setting $N=n+1$ yields
\[
 \frac{\delta Z}{\ii\delta J(x)}
 =\sum_{n=0}^{\infty}\frac{\ii^n}{n!}
 \int\prod_{r=1}^n\dd^D y_r\,J(y_r)
 \langle0|T\phi(x)\phi(y_1)\cdots\phi(y_n)|0\rangle.
\]
For example, the term linear in the remaining source is
\[
 \ii\int\dd^D y\,J(y)\left[
 \theta(x^0-y^0)\langle0|\phi(x)\phi(y)|0\rangle
 +\theta(y^0-x^0)\langle0|\phi(y)\phi(x)|0\rangle\right].
\]
Thus the newly exposed point $x$ participates in the same ordering operation
as $y$.  Taking it outside $T$ would discard the second term and change the
operator product whenever $y^0>x^0$.

The same counting gives, for any $r$,
\[
 \frac{\delta^r Z}
 {\ii\delta J(x_1)\cdots\ii\delta J(x_r)}
 =\left\langle0\left|T\phi(x_1)\cdots\phi(x_r)
 \exp\!\left(\ii\int\dd^D y\,J(y)\phi(y)\right)\right|0\right\rangle.
\]
If $V'(u)=\sum_r c_r u^r$, multiplication by $c_r$ followed by this identity
gives
\[
 V'\!\left(\frac{\delta}{\ii\delta J(x)}\right)Z[J]
 =\left\langle0\left|T V'(\phi(x))
 \exp\!\left(\ii\int\dd^D y\,J(y)\phi(y)\right)\right|0\right\rangle.
\]
This establishes the insertion in Eq.~\eqref{eq:3.4}.  Time derivatives of
the resulting expression require the contact terms developed immediately
after that equation.

% BANKS-SOURCE: pdf=31 print=21 kind=solution id=I-CH03-003
\BanksImplicitSolution{I-CH03-003}
For the kernel used in Eq.~\eqref{eq:3.11}, differentiation with respect to a
source inserts a Fourier variable:
\[
 \frac{\partial Z}{\ii\partial J_B}
 =\int\dd^N\phi\,\phi^B\ee^{\ii S+\ii J\phi}.
\]
Multiplication by the source is converted by integration by parts:
\begin{align*}
 J_A Z
 &=-\ii\int\dd^N\phi\,\ee^{\ii S}
        \frac{\partial}{\partial\phi^A}\ee^{\ii J\phi}\\
 &=-\int\dd^N\phi\,
       \frac{\partial S}{\partial\phi^A}\ee^{\ii S+\ii J\phi}.
\end{align*}
The vanishing boundary term can be enforced by the same small damping that
defines the oscillatory integral.

Use the convention implicit in Eqs.~\eqref{eq:3.12}--\eqref{eq:3.14}, in
which $K$ is the regulated quadratic kernel.  With the $+\ii J\phi$ kernel, the
Schwinger--Dyson equation consistent with the classical field equation is
\[
 K_{AB}\frac{\partial Z}{\ii\partial J_B}
 =\left[V'\!\left(\frac{\partial}{\ii\partial J_A}\right)-J_A\right]Z.
\]
Using the two transform rules reduces it to
\[
 \int\dd^N\phi\left(K_{AB}\phi^B
 -\frac{\partial V}{\partial\phi^A}
 -\frac{\partial S}{\partial\phi^A}\right)
 \ee^{\ii S+\ii J\phi}=0.
\]
Since this holds for every $J$, Fourier uniqueness gives
\[
 \frac{\partial S}{\partial\phi^A}
 =K_{AB}\phi^B-\frac{\partial V}{\partial\phi^A}.
\]
Symmetry of $K$ makes the right-hand side a gradient, and hence
\[
 S(\phi)=\frac12\phi^AK_{AB}\phi^B-V(\phi)+C.
\]
Differentiation recovers the preceding equation.  The constant multiplies
$Z$ by $\ee^{\ii C}$, which the homogeneous Schwinger--Dyson equations leave
undetermined.

Equation~\eqref{eq:3.9}, with the mostly-plus scalar equation
$\Box\phi=V'(\phi)$, gives the displayed $V'-J_A$ combination.  The historical
source version of Eq.~\eqref{eq:3.10} instead printed $+J_A$.  Combining that
printed sign with the kernel in
Eq.~\eqref{eq:3.11} would give $S_{,A}=V_{,A}-K_{AB}\phi^B$, opposite to
Eq.~\eqref{eq:3.12}.  Thus the latter equation fixes the source term in
Eq.~\eqref{eq:3.10} to $-J_A$.  Choosing the kernel $\ee^{-\ii J\phi}$
reverses both transform rules and requires the corresponding source-sign
change.

% BANKS-SOURCE: pdf=32 print=22 kind=solution id=I-CH03-004
\BanksImplicitSolution{I-CH03-004}
Set $a=\ii K^{-1}J$.  A direct expansion gives
\[
 -\frac12y^TKy+\ii y^TJ
 =-\frac12(y-a)^TK(y-a)-\frac12J^TK^{-1}J.
\]
The integrand is entire in every $y^i$.  Positive definiteness of
$\operatorname{Re}K$ supplies Gaussian decay on the original contour and on
the finite interpolation to $\mathbb R^n-a$, so Cauchy's theorem permits the
shift.  It follows that
\[
 I(J)=I(0)\exp\!\left(-\frac12J^TK^{-1}J\right).
\]

For a real positive-definite $K$, an orthogonal diagonalization and $n$
ordinary one-dimensional integrals give
\[
 I(0)=(2\pi)^{n/2}(\det K)^{-1/2}.
\]
Both sides are analytic throughout the connected domain
$\operatorname{Re}K>0$.  This fixes the general answer as
\[
 \boxed{I(J)=(2\pi)^{n/2}(\det K)^{-1/2}
 \exp\!\left(-\frac12J^TK^{-1}J\right)},
 \qquad
 \boxed{\frac{I(J)}{I(0)}=\ee^{-J^TK^{-1}J/2}}.
\]
The square root is the analytic branch reached from positive real $K$.
Replacing it by $|\det K|^{1/2}$ would lose the phase for a complex matrix.

For a real nonsingular quadratic form $A$, the Fresnel integral is obtained
from $K_\varepsilon=\varepsilon\mathbf1-\ii A$ with
$\varepsilon\downarrow0$.  If $A$ has $n_+$ positive and $n_-$ negative
eigenvalues, then
\[
 \int\dd^ny\,\ee^{\ii y^TAy/2}
 =(2\pi)^{n/2}|\det A|^{-1/2}
   \ee^{\ii\pi(n_+-n_-)/4},
\]
with the phase fixed by that limiting prescription.  The source term is
obtained from the same completed square.

At $n=1$, $J=0$, the defining integral is
$\int\dd y\,\ee^{-ky^2/2}=\sqrt{2\pi/k}$.  Equation~\eqref{eq:3.16}, as
printed, instead supplies $\sqrt{\pi/(2k)}$.  Its overall constant is smaller
by a factor of $2$ per integration variable, and its absolute determinant
does not record the complex phase.  Both discrepancies cancel from
$I(J)/I(0)$, which is the normalized functional integral used in the next
paragraph of the text.

% BANKS-SOURCE: pdf=32 print=22 kind=solution id=I-CH03-005
\BanksImplicitSolution{I-CH03-005}
Put $t=x^0-y^0$ and $\mathbf r=\mathbf x-\mathbf y$.  The annihilation part of
$\phi(x)$ can contract only with the creation part of $\phi(y)$.  Thus
\[
 D_F(t,\mathbf r)=
 \begin{cases}
 \displaystyle
 \int\frac{\dd^{D-1}\mathbf p}{(2\pi)^{D-1}\,2\omega_{\mathbf p}}
 \ee^{-\ii\omega_{\mathbf p}t+\ii\mathbf p\cdot\mathbf r},&t>0,\\[8pt]
 \displaystyle
 \int\frac{\dd^{D-1}\mathbf p}{(2\pi)^{D-1}\,2\omega_{\mathbf p}}
 \ee^{\ii\omega_{\mathbf p}t-\ii\mathbf p\cdot\mathbf r},&t<0.
 \end{cases}
\]
Changing $\mathbf p\mapsto-\mathbf p$ in the second line puts both cases in
the compact form
\[
 D_F(t,\mathbf r)=\int\frac{\dd^{D-1}\mathbf p}{(2\pi)^{D-1}\,2\omega_{\mathbf p}}
 \left[\theta(t)\ee^{-\ii\omega_{\mathbf p}t}
       +\theta(-t)\ee^{\ii\omega_{\mathbf p}t}\right]
 \ee^{\ii\mathbf p\cdot\mathbf r}.
\]

The denominator in
\[
 \int\frac{\dd p^0}{2\pi}\,
 \frac{\ii\ee^{-\ii p^0t}}
 {(p^0)^2-\omega_{\mathbf p}^2+\ii\varepsilon}
\]
has poles at $p^0=\omega_{\mathbf p}-\ii0$ and
$p^0=-\omega_{\mathbf p}+\ii0$.  For $t>0$ the exponential decays on the
lower semicircle and encloses the positive-energy pole.  For $t<0$ the upper
semicircle encloses the negative-energy pole.  The residues reproduce the
two preceding terms, proving
\[
 D_F(x-y)=\int\frac{\dd^Dp}{(2\pi)^D}
 \frac{-\ii\ee^{+\ii p\mathbin{\cdot}(x-y)}}
 {p^2+m^2-\ii\varepsilon}.
\]

Away from $t=0$, each Fourier mode solves the homogeneous Klein--Gordon
equation.  Its first derivative has the jump
\begin{align*}
 \partial_tD_F(0^+,\mathbf r)-\partial_tD_F(0^-,\mathbf r)
 &=-\ii\int\frac{\dd^{D-1}\mathbf p}{(2\pi)^{D-1}}\,
 \ee^{\ii\mathbf p\cdot\mathbf r}\\
 &=-\ii\delta^{D-1}(\mathbf r).
\end{align*}
Consequently, in the mostly-plus convention used by Banks,
\begin{align*}
 -\ii(\Box-m^2+\ii\varepsilon)D_F(x-y)&=\delta^D(x-y),\\
 (\Box-m^2)D_F(x-y)&=\ii\delta^D(x-y)
 \quad(\varepsilon\to0^+).
\end{align*}
For $t>0$ the positive-frequency part propagates from $y$ to $x$; for $t<0$
it propagates from $x$ to $y$.  This operator boundary condition is precisely
the placement of the two poles produced by $+\ii\varepsilon$.

% BANKS-SOURCE: pdf=33 print=23 kind=solution id=I-CH03-006
\BanksImplicitSolution{I-CH03-006}
Write $x_r=(t_r,\mathbf x_r)$ and let
$\{|a\rangle\}$ be a complete set with $H|a\rangle=E_a|a\rangle$ and
$E_a\geq0$.  Translation in time gives, in the ordering
$t_1>\cdots>t_n$,
\begin{align*}
 G_n(x_1,\ldots,x_n)
 ={}&\sum_{a_1,\ldots,a_{n-1}}
 \langle0|\phi(0,\mathbf x_1)|a_1\rangle
 \langle a_1|\phi(0,\mathbf x_2)|a_2\rangle\cdots\\
 &\quad\times
 \langle a_{n-1}|\phi(0,\mathbf x_n)|0\rangle
 \exp\!\left[-\ii\sum_{r=1}^{n-1}E_{a_r}(t_r-t_{r+1})\right].
\end{align*}
Continuous parts of the spectrum replace the sums by spectral integrals.
Equivalently, the products of matrix elements define a complex measure
$\dd\rho(E_1,\ldots,E_{n-1})$ supported on $E_r\geq0$, and
\[
 G_n=\int\dd\rho(E_1,\ldots,E_{n-1})
 \exp\!\left(-\ii\sum_rE_r z_r\right),
 \qquad z_r=t_r-t_{r+1}.
\]

Let $z_r=s_r-\ii\tau_r$ with every $\tau_r>0$.  A sufficient hypothesis is
\[
 \int |\dd\rho(E)|\,(1+E_1+\cdots+E_{n-1})^N
 \ee^{-\sum_r\tau_r E_r}<\infty
\]
for every nonnegative integer $N$ and every point in the tube.  Polynomially
bounded spectral measures satisfy this condition.  The exponential dominates
all derivatives with respect to the $z_r$, so dominated convergence proves
holomorphy.  At $z_r=-\ii\tau_r$ the kernel becomes
\[
 \ee^{-\sum_rE_r\tau_r}.
\]
The fields remain in the order $1,\ldots,n$ because $\tau_r>0$ preserves the
ordered chamber.  A different real-time ordering is the boundary value of a
different chamber.  The pieces agree where locality permits their analytic
continuation to meet.

For two free fields the operator result from the previous exercise continues
to
\[
 D_F(-\ii\tau,\mathbf r)
 =\int\frac{\dd^{D-1}\mathbf p}{(2\pi)^{D-1}\,2\omega_{\mathbf p}}
 \ee^{-\omega_{\mathbf p}|\tau|+\ii\mathbf p\cdot\mathbf r}.
\]
Starting from the $D$-momentum representation, rotate $p^0$ to
$\ii p_E^0$ together with $t=-\ii\tau$.  The Feynman poles lie on the safe
sides of the rotating quadrants.  The integrand is $O(|p^0|^{-2})$, and the
continued exponential damps the arc for separated points; a test-function
smearing supplies the same conclusion at coincident points.  Accounting for
$\dd p^0=\ii\dd p_E^0$ gives
\[
 D_E(\tau,\mathbf r)=\int\frac{\dd p_E^0\,\dd^{D-1}\mathbf p}{(2\pi)^D}
 \frac{\ee^{-\ii p_E^0\tau+\ii\mathbf p\mathbin{\cdot}\mathbf r}}
 {(p_E^0)^2+\mathbf p^2+m^2},
\]
which is Eq.~\eqref{eq:3.21} after the spatial change of variable
$\mathbf p_E=-\mathbf p$ in the Euclidean Fourier integral.

% BANKS-SOURCE: pdf=33 print=23 kind=solution id=I-CH03-007
\BanksImplicitSolution{I-CH03-007}
Smear the field only in space,
\[
 \Phi_f(t)=\int\dd^{D-1}\mathbf x\,f(\mathbf x)\phi(t,\mathbf x),
 \qquad f\in\mathcal S(\mathbb R^{D-1}).
\]
The free two-point function is
\[
 \langle0|\Phi_f(t)\Phi_f(0)|0\rangle
 =\int\frac{\dd^{D-1}\mathbf p}{(2\pi)^{D-1}\,2\omega_{\mathbf p}}
 |\widetilde f(\mathbf p)|^2\ee^{-\ii\omega_{\mathbf p}t}
 =\int_m^\infty\dd E\,\rho_f(E)\ee^{-\ii Et},
\]
where, with $\kappa_E=\sqrt{E^2-m^2}$,
\[
 \rho_f(E)=\frac{\theta(E-m)\kappa_E^{D-3}}{2(2\pi)^{D-1}}
 \int\dd\Omega_{D-2}\,|\widetilde f(\kappa_E\widehat{\mathbf p})|^2.
\]
At finite spatial volume the same object is a sum of delta functions.  In
infinite volume it has a $(D-3)/2$-power threshold (a square-root threshold when
$D=4$).  A Schwartz $f$ makes it decrease faster than every power as $E\to\infty$.
Point evaluation
corresponds formally to constant $\widetilde f$ and, in $D=4$, gives
$\rho(E)=\kappa_E/(4\pi^2)$, which grows.  Hence the unsmeared free spectral object
already has a threshold and is a measure rather than a globally analytic,
decreasing function.

At a fixed perturbative order, the discontinuity of a graph has the schematic
Cutkosky form
\[
 \operatorname{Disc}\mathcal G(P)
 =\sum_X\int\dd\Pi_X(P)\,
 \mathcal M_L(P\to X)\mathcal M_R(P\to X)^*,
\]
where $\dd\Pi_X$ contains on-shell delta functions and has support above the
multi-particle threshold for $X$.  Phase-space boundaries produce further
square roots and logarithms.  Renormalized fixed-order graphs are tempered:
power counting bounds their high-energy growth by powers and logarithms.  Full
space-time smearing multiplies them by Schwartz Fourier transforms and gives
rapid decrease.  Spatial smearing alone can leave polynomial growth in energy,
which the Euclidean Laplace factor controls.

The precise order-by-order statement is therefore distributional.  Spectral
distributions are tempered and supported at nonnegative energy.  After full
test-function smearing they decrease rapidly; under spatial smearing alone a
polynomial bound is enough.  Their Laplace transforms
\[
 \int_0^\infty\dd\rho_f(E)\,\ee^{-\ii Ez},
 \qquad \operatorname{Im}z<0,
\]
are analytic because $\ee^{E\operatorname{Im}z}$ controls every polynomial
bound.  The Euclidean value $z=-\ii\tau$ converges for every $\tau>0$.
Analyticity belongs to this complex-time transform and, away from cuts, to
momentum-space amplitudes.  The literal claim that the unsmeared spectral
density itself is analytic and decreasing fails at zeroth order; the
distributional formulation supplies exactly the continuation used by Banks.

% BANKS-SOURCE: pdf=35 print=25 kind=solution id=I-CH03-008
\BanksImplicitSolution{I-CH03-008}
Introduce normalized Gaussian expectation values
\[
 \langle F\rangle_0:=\frac{1}{I_0}
 \int[d\phi]\,\ee^{\ii S_0[\phi]}F[\phi],
 \qquad A:=\int\dd^D x\,V(\phi(x)),
\]
and abbreviate $a=-\ii g$.  Through second order,
\begin{align*}
 \frac{I[J,g]}{I_0}
 &=Z_0[J]+aN_1[J]+\frac{a^2}{2}N_2[J]+O(a^3),\\
 \frac{I[0,g]}{I_0}
 &=1+aD_1+\frac{a^2}{2}D_2+O(a^3),
\end{align*}
where
\[
 N_r[J]=\langle A^r\ee^{\ii\int\dd^D x\,J(x)\phi(x)}\rangle_0,
 \qquad D_r=\langle A^r\rangle_0.
\]
The factor $a=-\ii g$ follows by expanding
$\ee^{\ii(S_0-gA)}$; Eq.~\eqref{eq:3.26} already displays this
$(-\ii g)^n$ factor.
Every coefficient is visibly a Gaussian integral divided by $I_0$.  Expanding
the reciprocal denominator gives
\begin{align*}
 Z[J,g]={}&Z_0[J]+a\bigl(N_1-Z_0D_1\bigr)\\
 &+a^2\left[\frac12N_2-N_1D_1
 +Z_0\left(D_1^2-\frac12D_2\right)\right]+O(a^3).
\end{align*}
The subtractions remove the first- and second-order collections that contain
a component made solely from interaction vertices.

The all-orders statement follows from the exponential formula for graphs.
The logarithm of a Gaussian expectation is the sum of its connected Wick
graphs.  Thus
\begin{align*}
 \frac{I[J,g]}{I_0}
 &=\exp\!\left(C_{\rm vac}(g)+C_{\rm src}[J,g]\right),\\
 \frac{I[0,g]}{I_0}&=\exp\!\left(C_{\rm vac}(g)\right).
\end{align*}
Here $C_{\rm vac}$ contains connected vacuum graphs, while every term in
$C_{\rm src}$ meets at least one source insertion.  Division cancels
$C_{\rm vac}$ and every disconnected product containing it.  At $g=0$ this
leaves the normalized Gaussian generator
\[
 Z_0[J]=\langle\ee^{\ii\int\dd^D x\,J(x)\phi(x)}\rangle_0
 =\exp\!\left[-\frac12\int\dd^D x\,\dd^D y\,
 J(x)D_F(x-y)J(y)\right].
\]
Setting $J=0$ in the defining ratio gives $Z[0,g]=1$ exactly, which checks the
vacuum-bubble cancellation at every order.

% BANKS-SOURCE: pdf=35 print=25 kind=solution id=I-CH03-009
\BanksImplicitSolution{I-CH03-009}
Take the definite normalization
\[
 \mathcal L_I=-\frac{g}{2!2!}\phi^2
 \partial_\mu\phi\,\partial^\mu\phi.
\]
At a vertex $y$, temporarily distinguish the fields as
\[
 \phi_1(y)\phi_2(y)\,
 \partial_\mu\phi_3(y)\partial^\mu\phi_4(y).
\]
For four labeled external points, choose an unordered pair $\{a,b\}$ for the
differentiated slots.  The two differentiated fields can be assigned in
$2!$ ways, and the undifferentiated fields in another $2!$ ways.  This factor
cancels the $2!2!$ in the interaction.  The one-vertex contribution is
\[
 -\ii g\int\dd^D y\sum_{1\leq a<b\leq4}
 \partial_\mu^yD_F(y-x_a)\,
 \partial_y^\mu D_F(y-x_b)
 \prod_{c\ne a,b}D_F(y-x_c).
\]
The same rule applies to an internal line: a derivative attached to its slot
acts on the $y$ endpoint of that line's propagator.  For a line joining $y$
to $z$, this is $\partial_yD_F(y-z)$.  A derivative belonging to the $z$
vertex acts as $\partial_z$ and carries the opposite routed momentum.

Use
\[
 D_F(y-x)=\int\frac{\dd^D p}{(2\pi)^D}
 \widetilde D_F(p)\ee^{+\ii p\mathbin{\cdot}(y-x)}.
\]
Each differentiated slot contributes $+\ii p_\mu$.  Combining the two such
factors with the insertion $-\ii g$ gives the all-incoming vertex
\[
 \boxed{\mathcal V_4(p_1,p_2,p_3,p_4)
 =\ii g\sum_{a<b}p_a\mathbin{\cdot}p_b},
 \qquad (2\pi)^D\delta^D\!\left(\sum_a p_a\right).
\]
The sum makes permutation symmetry manifest.  Momentum conservation gives
\[
 \sum_{a<b}p_a\mathbin{\cdot}p_b
 =\frac12\left[\left(\sum_a p_a\right)^2-\sum_a p_a^2\right]
 =-\frac12\sum_a p_a^2.
\]
For four on-shell equal-mass legs, with the surrounding mostly-plus convention
$p_a^2=-m^2$, the vertex becomes $+2\ii g m^2$.  The
identity
\[
 \phi^2(\partial\phi)^2
 =\frac13\partial_\mu(\phi^3\partial^\mu\phi)
  -\frac13\phi^3\Box\phi
\]
shows why it vanishes for massless external legs.  For $m\ne0$, the free
equation of motion trades it for $-m^2\phi^4/3$ modulo a total derivative,
which reproduces the nonzero massive vertex.  These are the covariant
functional-integral rules; a canonical Dyson expansion with time derivatives
generates contact terms that reconstruct the same covariant result.

% BANKS-SOURCE: pdf=35 print=25 kind=solution id=I-CH03-010
\BanksImplicitSolution{I-CH03-010}
For real scalars with canonical kinetic terms, fix
\[
 \mathcal L_0=-\frac12\partial_\mu\phi_a\partial^\mu\phi_a
 -\frac12\phi_aM^2_{ab}\phi_b,
 \qquad
 \mathcal L_I=-\sum_{r\geq3}\frac{1}{r!}
 \lambda_{a_1\ldots a_r}\phi_{a_1}\cdots\phi_{a_r},
\]
where $M^2$ is real symmetric and each $\lambda$ is totally symmetric.  An
unsymmetrized coefficient is replaced by its complete symmetrization.
The quadratic momentum-space action kernel is
\[
 K_{ab}(p)=-p^2\delta_{ab}-M^2_{ab}+\ii\varepsilon\delta_{ab}.
\]
The Gaussian contraction of two fields is $\ii K^{-1}$:
\[
 \boxed{\widetilde D_F^{ab}(p)
 =-\ii\left[p^2\mathbf1+M^2-\ii\varepsilon\mathbf1\right]^{-1}_{ab}}.
\]
Consequently a line has an index at each end, and the two indices are joined
by $D_F^{ab}$.  Wick's theorem permits precisely the pairings encoded by this
covariance matrix.  In a basis where $M^2$ is diagonal,
\[
 \widetilde D_F^{\alpha\beta}(p)
 =\frac{-\ii\delta^{\alpha\beta}}
 {p^2+m_\alpha^2-\ii\varepsilon},
\]
so a free contraction preserves the mass-eigenstate species.

Expanding $\ee^{\ii\int\dd^D x\,\mathcal L_I}$ and contracting the $r!$ identical
slots with labeled incident lines cancels the explicit $1/r!$.  The vertex is
\[
 \boxed{-\ii\lambda_{a_1\ldots a_r}},
\]
together with momentum conservation.  An external line carries its specified
species index.  Every internal endpoint index is summed, as in
\[
 \lambda_{abc}D_F^{aa'}(p)D_F^{bb'}(q)
 \lambda_{a'b'c'}D_F^{cc'}(r)
\]
for a representative contraction of two cubic tensors.  Zeros or invariant
tensors in $\lambda$ impose the allowed species combinations at a vertex.

Under an orthogonal field redefinition $\phi'=O\phi$,
\[
 M'^2=OM^2O^T,
 \qquad D'_F=OD_FO^T,
 \qquad
 \lambda'_{a_1\ldots a_r}
 =O_{a_1b_1}\cdots O_{a_rb_r}\lambda_{b_1\ldots b_r}.
\]
Every internal index joins two factors of $O$ and contracts to
$O^TO=\mathbf1$.  External indices rotate with the external states.  Hence
the numerical amplitude between correspondingly rotated states is basis
independent, and the diagonal-basis rules are the same tensor network written
in simpler coordinates.

% BANKS-SOURCE: pdf=35 print=25 kind=solution id=I-CH03-011
\BanksImplicitSolution{I-CH03-011}
Suppose $G$ has $n_\tau$ vertices of each interaction type $\tau$.  Let
$R_\tau$ be the group that permutes identical slots of that vertex; for a
scalar monomial $\phi^{k_\tau}/k_\tau!$, it is $S_{k_\tau}$.  Before any
contractions, the perturbation series treats the vertices and their slots as
labeled.  The relabeling group is
\[
 H=\prod_\tau\left(S_{n_\tau}\ltimes R_\tau^{n_\tau}\right),
 \qquad |H|=\prod_\tau n_\tau!|R_\tau|^{n_\tau}.
\]
It acts transitively on the labeled Wick contractions that give the same
abstract graph $G$.  The subgroup that leaves one such contraction unchanged
is precisely $\operatorname{Aut}(G)$, provided an automorphism fixes the
labeled external legs and preserves all incidences, valences, species, vertex
types, and derivative decorations.  Edge
permutations and the exchange of the two ends of a self-loop occur as elements
of this stabilizer.

Orbit--stabilizer therefore gives
\[
 N_G=\frac{|H|}{{|\operatorname{Aut}(G)|}}
\]
contractions in the orbit of $G$.  The exponential contributes $1/n_\tau!$
for each set of identical vertices, while the interaction normalization
contributes $1/|R_\tau|$ at each such vertex.  Multiplying these denominators
by $N_G$ leaves
\[
 \boxed{S_G=\frac{N_G}
 {\prod_\tau n_\tau!|R_\tau|^{n_\tau}}
 =\frac{1}{|\operatorname{Aut}(G)|}}.
\]

For the order-$\lambda$ two-point tadpole, choose a slot for the field at $x$
and then a slot for the field at $y$.  The remaining two slots contract with
each other, so
\[
 N_G=4\cdot3=12,
 \qquad S_G=\frac{12}{4!}=\frac12.
\]
The only automorphism fixing the two external legs flips the tadpole edge, so
$|\operatorname{Aut}(G)|=2$.

For two quartic vertices joined by four parallel lines, a contraction is a
bijection between their four slots.  There are $4!$ such bijections.  Hence
\[
 S_G=\frac{4!}{2!(4!)^2}=\frac1{48}.
\]
The automorphism group has $4!$ permutations of the parallel edges and an
exchange of the two vertices, giving order $48$.  This is the $1/48$ example
in Appendix~D.  The external $x$--$y$ contraction drawn there forms a separate
labeled component and supplies no further multiplicity.

% BANKS-SOURCE: pdf=36 print=26 kind=solution id=I-CH03-012
\BanksImplicitSolution{I-CH03-012}
Adopt
\[
 \mathcal L_I=-g\sum_k\frac{v_k}{k!}\phi^k.
\]
Expansion of $\ee^{\ii S_I}=\ee^{\ii\int\dd^D x\,\mathcal L_I}$ then assigns
\[
 \text{vertex of valence }k:\quad -\ii g v_k,
 \qquad
 \text{internal line}:\quad
 \widetilde D_F(p)=\frac{-\ii}{p^2+m^2-\ii\varepsilon}.
\]
The $k!$ contractions of labeled incident lines cancel the $1/k!$ in the
Lagrangian.  A Green function also has one propagator from each external point
to the first vertex it meets; those factors are omitted from an amputated
quantity.

Insert the Fourier representation of every propagator into a position-space
graph.  Integration over a vertex position $y_v$ gives
\[
 \int\dd^D y_v\,\ee^{+\ii y_v\mathbin{\cdot}\sum_{e\ni v}k_{ev}}
 =(2\pi)^D\delta^D\!\left(\sum_{e\ni v}k_{ev}\right),
\]
where all $k_{ev}$ point into the vertex.  For a connected graph, one of the
$V$ vertex constraints becomes the overall conservation factor.  The nearby
source line prints $2\pi^4\delta^4$ in its four-dimensional presentation;
the D-dimensional Fourier transformation gives the parenthesized factor
$(2\pi)^D\delta^D$ displayed here.
The remaining $V-1$ delta functions eliminate $V-1$ of the $I$ internal
momenta.  Therefore
\[
 L=I-(V-1)=I-V+1
\]
independent loop momenta remain.  The resulting amputated graph is
\begin{align*}
 \mathcal G_G={}&(2\pi)^D\delta^D\!\left(\sum_a p_a\right)S_G
 \prod_{v=1}^{V}(-\ii g v_{k_v})
 \int\prod_{r=1}^{L}\frac{\dd^D\ell_r}{(2\pi)^D}
 \prod_{e=1}^{I}
 \frac{-\ii}{k_e(\ell,p)^2+m^2-\ii\varepsilon}.
\end{align*}

With the convention $\ee^{+\ii k\mathbin{\cdot}x}$, a derivative acting on a field at a
vertex contributes $+\ii k_\mu$ for that field's incoming momentum.  Several
derivatives produce the corresponding tensor polynomial.  Momentum
conservation is imposed only after each derivative has been assigned to its
own field slot.

Under $t=-\ii\tau$, this scalar action becomes
\[
 S_E=\int\dd^D x_E\left[\frac12(\partial_E\phi)^2
 +\frac12m^2\phi^2+g\sum_k\frac{v_k}{k!}\phi^k\right],
 \qquad \ee^{\ii S_M}=\ee^{-S_E}.
\]
Thus
\[
 \text{Euclidean line}:\quad\frac{1}{p_E^2+m^2},
 \qquad
 \text{Euclidean vertex}:\quad-gv_k,
 \qquad
 \int\frac{\dd^D\ell_E}{(2\pi)^D}
\]
for each loop.  The Lorentzian coefficient shown above follows from
$\mathcal L_I=-g\sum_kv_k\phi^k/k!$ and $\ee^{\ii S_M}$, with the same
coupling convention used in the Euclidean rules.  The displayed signs follow
directly from $\ee^{\ii S_M}$ and $\ee^{-S_E}$.

% BANKS-SOURCE: pdf=36 print=26 kind=solution id=I-CH03-013
\BanksImplicitSolution{I-CH03-013}
Use $\mathcal L_I=-\lambda\phi^4/4!$, set
$D_{ab}=D_F(x_a-x_b)$, and retain all external propagators.  The phases below
follow this choice.  Appendix~D uses the opposite coupling sign; its graph
counts are unchanged.  The one-vertex two-point tadpole is
\[
 G^{(1)}_2(x,y)=\frac{-\ii\lambda}{2}
 \int\dd^D z\,D_F(x-z)D_F(0)D_F(z-y).
\]
It has one vertex, one internal self-loop, one loop momentum, and two external
legs.  The count $4\cdot3/4!=1/2$ gives its symmetry factor.  Its superficial
degree of divergence is $D-2$, which is $2$ when $D=4$.

The connected four-point contact graph is
\[
 G^{(1)}_4(x_1,x_2,x_3,x_4)
 =-\ii\lambda\int\dd^D z\prod_{a=1}^4D_F(x_a-z).
\]
All $4!$ assignments of external points cancel $1/4!$.  The graph has four
external legs, no internal line, no loop, and symmetry factor $1$.

For the two-point sunset, take $x$ into $z$, $y$ into $w$, and join $z$ to
$w$ with three lines:
\[
 G^{(2)}_{2,\mathrm{sun}}(x,y)
 =\frac{(-\ii\lambda)^2}{6}
 \int\dd^D z\,\dd^D w\,
 D_F(x-z)D_F(z-w)^3D_F(w-y).
\]
Here $V=2$, $I=3$, and $L=2$.  Directly,
\[
 \frac{2\cdot4\cdot4\cdot3!}{2!(4!)^2}=\frac16.
\]
The leading factor $2$ chooses which vertex receives $x$; the three remaining
lines can then be paired in $3!$ ways.  Its overall superficial degree is
$DL-2I=2D-6$, which is $2$ when $D=4$.

Choose the $s$-channel partition $(x_1,x_2)|(x_3,x_4)$ for the four-point
bubble.  Its position-space expression is
\[
\begin{aligned}
 G^{(2)}_{4,s}
 &=\frac{(-\ii\lambda)^2}{2}
   \int\dd^D z\,\dd^D w\\
 &\quad{}\times D_F(x_1-z)D_F(x_2-z)D_F(z-w)^2\\
 &\quad{}\times D_F(x_3-w)D_F(x_4-w).
\end{aligned}
\]
It has $V=2$, $I=2$, $L=1$, and symmetry factor $1/2$ from exchange of the
two parallel internal lines.  With $P=p_1+p_2$ and all external momenta
incoming, the amputated integral is
\[
 \mathcal G^{\mathrm{amp}}_{4,s}
 =\frac{(-\ii\lambda)^2}{2}
 \int\frac{\dd^D\ell}{(2\pi)^D}
 \frac{-\ii}{\ell^2+m^2-\ii\varepsilon}
 \frac{-\ii}{(\ell+P)^2+m^2-\ii\varepsilon},
\]
along with $(2\pi)^D\delta^D(\sum_ap_a)$.  The $t$ and $u$ channels are the
other two pairings of the four external legs.  Since $DL-2I=D-4$, each bubble
is logarithmically divergent when $D=4$.

Appendix~D also displays the vacuum component made from two quartic vertices
and four parallel lines.  Its factor
\(4!/[2!(4!)^2]=1/48\) agrees with the automorphism count in the previous
solution.  These examples exhibit every part of the setup before any loop
integration is attempted.

% BANKS-SOURCE: pdf=39 print=29 kind=solution id=I-CH03-014
\BanksImplicitSolution{I-CH03-014}
Differentiating $J=-\delta\Gamma/\delta\phi$ and
$\phi=\delta W/\delta J$ gives the kernel identity
\[
 \int\dd^D z\,W_2(x,z)[-\Gamma_2(z,y)]=\delta^D(x-y).
\]
Consequently
\[
 \boxed{W_2=-\Gamma_2^{-1}}.
\]
For any invertible kernel $K$,
\(\delta K^{-1}=-K^{-1}(\delta K)K^{-1}\).  Applying this to the boxed
identity yields
\begin{align*}
 W_3(x,y,z)
 &=\int\dd^D a\,\dd^D b\,
 W_2(x,a)W_2(y,b)\frac{\delta\Gamma_2(a,b)}{\delta J(z)},\\
 \frac{\delta\Gamma_2(a,b)}{\delta J(z)}
 &=\int\dd^D c\,\Gamma_3(a,b,c)W_2(c,z).
\end{align*}
Combining them gives
\[
 W_3(x,y,z)=\int\dd^D a\,\dd^D b\,\dd^D c\,
 W_2(x,a)W_2(y,b)W_2(z,c)\Gamma_3(a,b,c),
\]
which is Eq.~\eqref{eq:3.38}.

For four external points, define
\begin{align*}
 T_{12|34}:={}&\int\dd^D a\,\dd^D b\,\dd^D c\,\dd^D d\,
                 \dd^D u\,\dd^D v\,
 W_2(x_1,a)W_2(x_2,b)\Gamma_3(a,b,u)\\
 &\qquad\times W_2(u,v)\Gamma_3(v,c,d)
 W_2(x_3,c)W_2(x_4,d).
\end{align*}
One more differentiation gives
\begin{align*}
 W_4(x_1,x_2,x_3,x_4)
 ={}&\int\prod_{r=1}^4\dd^D a_r\,
 \left[\prod_{r=1}^4W_2(x_r,a_r)\right]
 \Gamma_4(a_1,a_2,a_3,a_4)\\
 &+T_{12|34}+T_{13|24}+T_{14|23}.
\end{align*}
The first term comes from differentiating $\Gamma_3$; differentiating one of
the three propagators produces the three exchange channels.

Assume the tree formula holds for $W_n$ with its external points labeled.
When $\delta/\delta J(x_{n+1})$ hits a vertex $\Gamma_k$, the chain rule
creates a vertex $\Gamma_{k+1}$ and attaches the new point through one
$W_2$.  When it hits an edge $W_2$, the three-point identity replaces that
edge by two edges meeting a new $\Gamma_3$, again with an external $W_2$ for
$x_{n+1}$.  Both operations preserve connectedness and create no cycle.

Conversely, remove the leaf labeled $x_{n+1}$ from any resulting tree.  If
its adjacent vertex has valence at least four, the inverse operation lowers
that vertex's valence and identifies a unique vertex differentiation.  If the
adjacent vertex is trivalent, removing it and joining its other two branches
identifies a unique differentiated edge.  This inverse establishes existence
and uniqueness, completing the induction.  In the free theory every
$\Gamma_{k\geq3}$ vanishes, and the formula correctly gives
$W_{n\geq3}=0$.

% BANKS-SOURCE: pdf=43 print=33 kind=solution id=I-CH03-015
\BanksImplicitSolution{I-CH03-015}
Translation covariance reads
\(\Phi(x)=\ee^{-\ii P\mathbin{\cdot}x}\Phi(0)\ee^{+\ii P\mathbin{\cdot}x}\).  Since
$P^\mu|0\rangle=0$ and $P^\mu|\mathbf p\rangle=p^\mu|\mathbf p\rangle$,
\[
 \langle0|\Phi(x)|\mathbf p\rangle
 =\ee^{+\ii p\mathbin{\cdot}x}F(\mathbf p),
 \qquad F(\mathbf p):=\langle0|\Phi(0)|\mathbf p\rangle.
\]
The origin is fixed by every Lorentz transformation, so scalar covariance,
vacuum invariance, and the stated transformation of the one-particle state
give
\[
 F(\boldsymbol{\Lambda p})
 =\sqrt{\frac{\omega_{\mathbf p}}{\omega_{\boldsymbol{\Lambda p}}}}\,F(\mathbf p).
\]
It follows that $\sqrt{\omega_{\mathbf p}}F(\mathbf p)$ is invariant on the positive
mass shell.  The proper orthochronous Lorentz group acts transitively there,
so this quantity is a constant.  Define
\[
 Z:=2(2\pi)^{D-1}\omega_{\mathbf p}|F(\mathbf p)|^2.
\]
It is independent of $p$ and positive when the interpolating field couples to
the particle.  A common phase choice for the one-particle states makes $F$
positive at rest.  The result is
\[
 \boxed{\langle0|\Phi(x)|\mathbf p\rangle
 =\sqrt{\frac{Z}{(2\pi)^{D-1}\,2\omega_{\mathbf p}}}\ee^{+\ii p\mathbin{\cdot}x}}.
\]

Covariantly normalized states are
\[
 |p\rangle_{\rm cov}=\sqrt{(2\pi)^{D-1}\,2\omega_{\mathbf p}}\,|\mathbf p\rangle,
 \qquad
 \langle q|p\rangle_{\rm cov}
 =(2\pi)^{D-1}2\omega_{\mathbf p}\delta^{D-1}(\mathbf q-\mathbf p).
\]
For them the same statement becomes
\(\langle0|\Phi(x)|p\rangle_{\rm cov}
=\sqrt Z\ee^{+\ii p\mathbin{\cdot}x}\).

Insert the one-particle completeness relation into the time-ordered two-point
function.  Its one-particle contribution is $Z$ times the free propagator.
Under the LSZ assumption that this pole is isolated from the continuum, near
$p^2=-m^2$,
\[
 \widetilde G_2(p)=\frac{-\ii Z}{p^2+m^2-\ii\varepsilon}
 +\text{terms regular at }p^2=-m^2.
\]
Thus $Z$ defined by the matrix element is the pole residue, which provides the
LSZ check used in the surrounding text.

% BANKS-SOURCE: pdf=43 print=33 kind=solution id=I-CH03-016
\BanksImplicitSolution{I-CH03-016}
Use the same delta-function normalization for momentum states as in the
previous solution.  For a stable massive spin-one particle, covariance fixes
\[
 \langle0|A_\mu(x)|\mathbf p,\lambda\rangle
 =\sqrt{\frac{Z_A}{(2\pi)^{D-1}\,2\omega_{\mathbf p}}}
 \epsilon_\mu(p,\lambda)\ee^{+\ii p\mathbin{\cdot}x}.
\]
With the mostly-plus metric, physical polarizations obey
\[
 p^\mu\epsilon_\mu(p,\lambda)=0,
 \qquad
 \epsilon(p,\lambda)^*\mathbin{\cdot}\epsilon(p,\lambda')
 =+\delta_{\lambda\lambda'},
\]
and
\[
 \sum_{\lambda=1}^{D-1}\epsilon_\mu(p,\lambda)
 \epsilon_\nu(p,\lambda)^*
 =\Pi^{(1)}_{\mu\nu}(p)
 :=\eta_{\mu\nu}+\frac{p_\mu p_\nu}{m^2}.
\]
There are $D-1$ massive vector polarizations; this is three when $D=4$.
The isolated pole in the vector two-point function is therefore
\[
 \widetilde G_{\mu\nu}(p)
 =\frac{-\ii Z_A\Pi^{(1)}_{\mu\nu}(p)}
 {p^2+m^2-\ii\varepsilon}+\text{regular terms}.
\]
On a physical transverse leg, use
\[
 Z_A^{-1/2}D_{F,0}^{-1}(p)
 =\frac{p^2+m^2}{-\ii\sqrt{Z_A}}.
\]
Contract the remaining index with $\epsilon^\mu$ for an incoming vector or
$\epsilon^{\mu *}$ for an outgoing vector.  A convention that amputates with
$p^2+m^2$ absorbs the displayed universal factor $1/(-\ii)$ into the definition
of the amputated vertex.

For a Dirac interpolating field,
the standard four-dimensional realization has $D=4$ and two spin states; in
other dimensions the same formulas use the corresponding $D$-dimensional
Clifford representation and spin multiplicity.
\begin{align*}
 \langle0|\Psi_\alpha(x)|\mathbf p,s\rangle
 &=\sqrt{\frac{Z_\psi}{(2\pi)^{D-1}\,2\omega_{\mathbf p}}}
 u_\alpha(p,s)\ee^{+\ii p\mathbin{\cdot}x},\\
 \langle0|\overline\Psi_\alpha(x)|\overline{\mathbf p},s\rangle
 &=\sqrt{\frac{Z_\psi}{(2\pi)^{D-1}\,2\omega_{\mathbf p}}}
 \overline v_\alpha(p,s)\ee^{+\ii p\mathbin{\cdot}x}.
\end{align*}
Choose
\begin{gather*}
 \{\gamma^\mu,\gamma^\nu\}=-2\eta^{\mu\nu},
 \qquad \slashed p=\gamma^\mu p_\mu,\\
 (\slashed p-m)u(p,s)=0,
 \qquad (\slashed p+m)v(p,s)=0,\\
 \overline u(p,s)u(p,s')=-2m\delta_{ss'},
 \qquad \overline v(p,s)v(p,s')=2m\delta_{ss'},\\
 u(p,s)^\dagger u(p,s')=v(p,s)^\dagger v(p,s')
 =2\omega_{\mathbf p}\delta_{ss'},\\
 \sum_su(p,s)\overline u(p,s)=-(\slashed p+m),
 \qquad \sum_sv(p,s)\overline v(p,s)=-(\slashed p-m).
\end{gather*}
The particle pole is
\[
 \widetilde S(p)=\frac{-\ii Z_\psi(\slashed p+m)}
 {p^2+m^2-\ii\varepsilon}+\text{regular terms}.
\]
Equivalently it is $\ii Z_\psi/(\slashed p-m+\ii0)$ near the pole.  Amputation
applies $Z_\psi^{-1/2}S_{F,0}^{-1}$ on the appropriate spinor index, where
\[
 S_{F,0}^{-1}(p)=\frac{\slashed p-m}{\ii}.
\]
After amputation, the external wave functions are
\[
 \begin{aligned}
 \text{incoming fermion}:&\quad Z_\psi^{-1/2}u(p,s),&
 \text{outgoing fermion}:&\quad Z_\psi^{-1/2}\overline u(p,s),\\
 \text{incoming anti-fermion}:&\quad Z_\psi^{-1/2}\overline v(p,s),&
 \text{outgoing anti-fermion}:&\quad Z_\psi^{-1/2}v(p,s).
 \end{aligned}
\]
The side on which $S_{F,0}^{-1}$ acts follows the oriented fermion line.  These
are the external-line conventions collected in Appendix~D.

For integer spin $s$, replace $\epsilon_\mu$ by a symmetric, transverse,
traceless tensor $\epsilon_{\mu_1\ldots\mu_s}(p,\lambda)$ normalized by
\[
\epsilon(p,\lambda)^*\mathbin{\cdot}\epsilon(p,\lambda')
 =\delta_{\lambda\lambda'}.
\]
Its polarization sum is the massive spin-$s$ projector
$\Pi^{(s)}(p)$, and the two-point pole is
$-\ii Z_s\Pi^{(s)}(p)/(p^2+m^2-\ii0)$.  For spin $s=r+\tfrac12$, use a
symmetric tensor--spinor $u_{\mu_1\ldots\mu_r}$ satisfying the Dirac equation,
transversality, and gamma tracelessness.  A compatible normalization is
\[
 \overline u_{\mu_1\ldots\mu_r}(p,\lambda)
 u^{\mu_1\ldots\mu_r}(p,\lambda')
 =-2m\delta_{\lambda\lambda'}.
\]
Its completeness relation is $(\slashed p+m)$ times the spin projector.
Each LSZ leg is obtained by applying the projected inverse free kinetic
operator, multiplying by $Z_s^{-1/2}$, and contracting with the appropriate
polarization tensor or tensor--spinor.

For a massless gauge particle, a covariant field contains gauge directions in
addition to physical states.  LSZ is applied to a gauge-fixed Green function
and then projected onto transverse physical helicities; BRST identities make
the result independent of the representative polarization.  This replaces
the massive polarization sum by its physical-helicity counterpart.

Setting $s=0$, $\Pi^{(0)}=1$, and using the scalar inverse pole recovers the
matrix element, residue, and $Z^{-1/2}$ reduction of the preceding solution.

\end{bankssolutions}
